\chapter{The Gaussian Information Form and the Interaction Family}
\label{ch:gaussian}

The preceding part fixed a normalized generative law, a recognition law that is permitted to be correlated, and one exact evidence identity relating them. Nothing in that construction chose a parametric form. This chapter makes the choice. It develops the Gaussian recognition family in information form, where the primitive coordinates are a vector and a matrix of natural parameters rather than a mean and a covariance, and it then isolates the particular subfamily of Gaussians on which the renormalization half of this document acts.

Two things are worth separating at the outset, because conflating them is the standing hazard of the whole subject. Writing a Gaussian in information form is a change of coordinates and costs nothing. Restricting the precision to the Laplacian-plus-self-terms shape of Section~\ref{sec:gauss-interaction-family} is a substantive restriction that costs a great deal, and it does not follow from the generative constructions of the preceding part. Section~\ref{sec:gauss-hypothesis} proves that it does not, exhibits the exact obstruction, and then states precisely the parameter condition under which the shape is recovered. The restriction is carried forward as a declared hypothesis, and every later result that uses it says so.

Throughout, the development is written for the state sector, with agents $i\in V=\{1,\dots,N\}$, state latents $k_i\in\R^{K}$, and stacked latent $Y$ of dimension $n:=NK$. The model sector, with latents $m_i\in\R^{d_m}$, is identical with $K$ replaced by $d_m$, and a partition mixing the two sectors is admissible wherever a partition is required. The design $D=\{c_a\}_{a=1}^{M}$ plays no role in the algebra below; either fix one design point, or read the index $i$ as ranging over agent-design pairs, and every statement stands unchanged.

\section{The recognition family in information form}
\label{sec:gauss-information-form}

Let $\nu$ be Lebesgue measure on $\R^{n}$. The Gaussian family in \emph{information form} is parameterized by a vector $h\in\R^{n}$ and a symmetric positive definite matrix $J\in\Sym_{++}^{n}$ through the density
\begin{equation}
q(y\given h,J)
=\exp\left(h^\top y-\tfrac12 y^\top Jy-\mathsf A(h,J)\right),
\qquad y\in\R^{n}.
\label{eq:gauss-density}
\end{equation}
The pair $(h,J)$ are the natural parameters, $J$ is the precision, and $\mathsf A$ is the log normalizer. The symbol $\mathsf A$ is reserved for the log normalizer and never for a self term; the self terms of Section~\ref{sec:gauss-interaction-family} are written $A_i$ with a subscript and are $K\times K$ matrices.

\paragraph{Proposition 6.1 (log normalizer, moments, convexity).}
For every $(h,J)\in\R^{n}\times\Sym_{++}^{n}$ the density in Equation~\eqref{eq:gauss-density} is normalizable, with
\begin{equation}
\mathsf A(h,J)
=\tfrac12 h^\top J^{-1}h-\tfrac12\log\det J+\tfrac n2\log 2\pi.
\label{eq:gauss-lognorm}
\end{equation}
Its differential is
\begin{equation}
d\mathsf A=\mu^\top dh-\tfrac12\Tr\left[\left(C+\mu\mu^\top\right)dJ\right],
\qquad
\mu=J^{-1}h,
\qquad
C=J^{-1},
\label{eq:gauss-moments}
\end{equation}
so the mean and covariance of $q(\cdot\given h,J)$ are $\mu$ and $C$, and the Hessian of $\mathsf A$ in $h$ at fixed $J$ is $C$. Regarded as a function of the natural parameter pair $(h,\Theta)$ with $\Theta=-\tfrac12J$, the log normalizer is strictly convex on the open convex domain $\R^{n}\times\{\Theta\prec0\}$. \status{ESTABLISHED}

\paragraph{Proof.}
Complete the square in the exponent,
\begin{equation}
h^\top y-\tfrac12 y^\top Jy
=-\tfrac12\left(y-J^{-1}h\right)^\top J\left(y-J^{-1}h\right)+\tfrac12 h^\top J^{-1}h,
\label{eq:gauss-complete-square}
\end{equation}
and integrate the Gaussian kernel, whose value is $(2\pi)^{n/2}(\det J)^{-1/2}$ because $J\succ0$. Exponentiating and taking logarithms gives Equation~\eqref{eq:gauss-lognorm}. Differentiating Equation~\eqref{eq:gauss-lognorm} termwise, $d(\tfrac12h^\top J^{-1}h)=h^\top J^{-1}dh-\tfrac12\Tr(J^{-1}hh^\top J^{-1}dJ)$ and $d(-\tfrac12\log\det J)=-\tfrac12\Tr(J^{-1}dJ)$, which is Equation~\eqref{eq:gauss-moments}. The identification of $\nabla_h\mathsf A$ with the mean and $\nabla_h^2\mathsf A$ with the covariance is the standard exponential-family cumulant property, valid here because the family is regular: the domain is open and differentiation under the integral is licensed on the interior of the natural domain \citep{Brown1986}. Convexity is the convexity of a log-Laplace transform, and strictness follows because the family is minimal, the statistics $(y,yy^\top)$ being affinely independent \citep{Brown1986,Wainwright2008}. Since $J\mapsto-\tfrac12J$ is affine, convexity in $(h,\Theta)$ transfers to convexity in $(h,J)$ on $\{J\succ0\}$. $\square$

Two consequences of Equation~\eqref{eq:gauss-moments} organize everything that follows. The first is that $(h,J)$ and $(\mu,C)$ are two coordinate systems on one family, related by the Legendre correspondence of a minimal exponential family, so no information is lost by working in either. The second is that the Fisher information of the family in the natural parameter $h$ is exactly $C=J^{-1}$, while the Fisher information in the mean parameter $\mu$ at fixed covariance is exactly $J$. The precision is therefore simultaneously a coupling matrix and a metric, and the renormalization chapters use both readings.

\section{What an off-diagonal precision block is}
\label{sec:gauss-conditional-reading}

Let $\mathfrak B$ be a finite partition of the $n$ coordinates into blocks, with block dimensions $n_b$ summing to $n$, and write $J_{bc}$ for the corresponding submatrices. The meaning of $J_{bc}$ for $b\neq c$ is fixed by the following statement, which is standard but is stated here in full because the rest of the document depends on reading it correctly.

\paragraph{Proposition 6.2 (conditional reading of the precision).}
Let $Y\sim\mathcal N(\mu,J^{-1})$ with $J\succ0$. For any block $b$,
\begin{equation}
\E\left[Y_b\given Y_{-b}=y_{-b}\right]
=\mu_b-J_{bb}^{-1}\sum_{c\neq b}J_{bc}\left(y_c-\mu_c\right),
\qquad
\operatorname{Cov}\left(Y_b\given Y_{-b}\right)=J_{bb}^{-1}.
\label{eq:gauss-conditional}
\end{equation}
For any two distinct blocks $b$ and $c$, the variables $Y_b$ and $Y_c$ are conditionally independent given all remaining coordinates if and only if $J_{bc}=0$. \status{ESTABLISHED}

\paragraph{Proof.}
Fix $y_{-b}$ and read the exponent of Equation~\eqref{eq:gauss-density} as a function of $y_b$ alone. It is $-\tfrac12y_b^\top J_{bb}y_b$ plus a linear term, so the conditional law is Gaussian with precision $J_{bb}$, and completing the square in $y_b$ around $\mu$ gives the displayed conditional mean. For the second claim, condition on the coordinates outside $b\cup c$ and repeat the argument: the conditional law of $(Y_b,Y_c)$ is Gaussian with precision $\begin{psmallmatrix}J_{bb}&J_{bc}\\J_{cb}&J_{cc}\end{psmallmatrix}$, and a Gaussian density factorizes into a function of $y_b$ times a function of $y_c$ if and only if the quadratic form has no cross term, that is if and only if $J_{bc}=0$. This is the Gaussian case of the pairwise Markov property \citep[Ch.~5]{Lauritzen1996}; the same statement organizes the Gaussian Markov random field literature \citep[Ch.~2]{RueHeld2005}. $\square$

Equation~\eqref{eq:gauss-conditional} also fixes the sign convention that Section~\ref{sec:gauss-interaction-family} will exploit. A negative semidefinite off-diagonal block $J_{bc}\preceq0$ makes the conditional mean of $Y_b$ move \emph{toward} $y_c$ as $y_c$ departs from its own mean, so negative off-diagonal precision is attractive coupling. The interaction family is the family in which every off-diagonal block is attractive in this sense.

\paragraph{Remark 6.3 (a precision block is not a covariance value).}
A zero off-diagonal precision block is a statement about conditional dependence and carries no implication about the marginal covariance. The tridiagonal example
\begin{equation}
J=\begin{pmatrix}2&-1&0\\-1&2&-1\\0&-1&2\end{pmatrix},
\qquad
J^{-1}=\frac14\begin{pmatrix}3&2&1\\2&4&2\\1&2&3\end{pmatrix},
\label{eq:gauss-precision-not-covariance}
\end{equation}
has $J_{13}=0$ and $(J^{-1})_{13}=\tfrac14\neq0$. In general the marginal covariance of a block is the inverse of a Schur complement,
\begin{equation}
\left(J^{-1}\right)_{bb}
=\left(J_{bb}-J_{b,-b}J_{-b,-b}^{-1}J_{-b,b}\right)^{-1},
\label{eq:gauss-marginal-block}
\end{equation}
so it records the effect of every path through the remaining coordinates, whereas $J_{bb}^{-1}$ records only the conditional dispersion of Equation~\eqref{eq:gauss-conditional} \citep{Zhang2005Schur}. The two coincide exactly when $J_{b,-b}=0$. \status{ESTABLISHED}

The distinction in Remark 6.3 is the reason this document works in information form. A restriction imposed on the recognition family is a restriction on conditional dependence, and it is stated cleanly only in the coordinates in which conditional dependence is a coordinate. The exact price of such a restriction is computed in the next chapter.

\section{The declared interaction subfamily}
\label{sec:gauss-interaction-family}

\paragraph{Definition 6.4 (interaction family).}
Let $\Lambda\in\Sym^{NK}$ be partitioned into $K\times K$ blocks indexed by agents. Say that $\Lambda$ belongs to the \emph{interaction family} $\mathcal I(V)$ when there exist matrices $A_i\in\PSD^{K}$ and $W_{ij}=W_{ji}\in\PSD^{K}$ for $i\neq j$ with
\begin{equation}
\Lambda_{ii}=A_i+\sum_{j\neq i}W_{ij},
\qquad
\Lambda_{ij}=-W_{ij}\quad(i\neq j).
\label{eq:gauss-interaction-family}
\end{equation}
The $A_i$ are the \emph{self terms}, the $W_{ij}$ the \emph{edge weights}, and $L:=\Lambda-\operatorname{blockdiag}(A_1,\dots,A_N)$ the \emph{Laplacian part}. Nothing is proved by this definition; it declares a type. \status{DEFINITION}

The edge weights are matrices, so $L$ is a matrix-weighted graph Laplacian rather than a scalar-weighted one. That object, and the ways in which it fails to behave like its scalar counterpart, are the subject of the matrix-weighted consensus literature \citep{Trinh2018MatrixWeighted}, and one of those failures appears immediately below.

\paragraph{Proposition 6.5 (energy, semidefiniteness, kernel).}
For $\Lambda\in\mathcal I(V)$ and any $y=(y_1,\dots,y_N)$,
\begin{equation}
y^\top\Lambda y
=\sum_{i\in V}y_i^\top A_iy_i
+\sum_{i<j}\left(y_i-y_j\right)^\top W_{ij}\left(y_i-y_j\right).
\label{eq:gauss-interaction-energy}
\end{equation}
Consequently $\Lambda\succeq0$ for every admissible choice of parameters; the consensus subspace $\mathbf 1\otimes\R^{K}$ lies in $\ker L$; and $\Lambda\succ0$ if and only if the only $y$ satisfying $A_iy_i=0$ for all $i$ together with $W_{ij}(y_i-y_j)=0$ for all $i<j$ is $y=0$. A necessary condition for $\Lambda\succ0$ is $\sum_iA_i\succ0$. \status{ESTABLISHED}

\paragraph{Proof.}
Expand the right-hand side of Equation~\eqref{eq:gauss-interaction-energy}. The self sum contributes $\sum_iy_i^\top A_iy_i$ to the diagonal blocks. Each unordered pair contributes $y_i^\top W_{ij}y_i+y_j^\top W_{ij}y_j-2y_i^\top W_{ij}y_j$, which places $W_{ij}$ on the $ii$ and $jj$ diagonal blocks and $-W_{ij}$ on the $ij$ and $ji$ off-diagonal blocks. Summing reproduces Equation~\eqref{eq:gauss-interaction-family}. Every term of Equation~\eqref{eq:gauss-interaction-energy} is nonnegative because $A_i\succeq0$ and $W_{ij}\succeq0$, so $\Lambda\succeq0$; and a sum of nonnegative terms vanishes only when each vanishes, which for a positive semidefinite matrix $A$ and vector $v$ means $v^\top Av=0$ if and only if $Av=0$. Taking $y=\mathbf 1\otimes v$ makes every difference vanish, so the consensus subspace lies in $\ker L$, and evaluating Equation~\eqref{eq:gauss-interaction-energy} there gives $v^\top(\sum_iA_i)v$, whence the necessary condition. $\square$

The kernel criterion in Proposition 6.5 is where the matrix-weighted setting departs from the scalar one. In the scalar case a connected interaction graph forces $\ker L$ to be exactly the consensus subspace. With matrix weights, a singular $W_{ij}$ leaves the difference $y_i-y_j$ free along $\ker W_{ij}$, so $\ker L$ can be strictly larger than the consensus subspace even on a connected graph, and connectivity alone is not a definiteness criterion \citep{Trinh2018MatrixWeighted}. The criterion stated in the proposition is exact and replaces it.

\paragraph{Proposition 6.6 (the family is thin).}
For $K\geq2$ and $N\geq2$, $\mathcal I(V)$ is contained in a proper linear subspace of $\Sym^{NK}$ of codimension at least $\binom N2\tfrac{K(K-1)}{2}$, hence has Lebesgue measure zero in $\Sym^{NK}$. \status{ESTABLISHED}

\paragraph{Proof.}
Symmetry of a matrix $\Lambda$ already forces $\Lambda_{ji}=\Lambda_{ij}^\top$ for every ordered pair, so this identity carries no information about the family. Membership additionally requires each off-diagonal block $\Lambda_{ij}$ to be symmetric \emph{as a $K\times K$ matrix}, since $W_{ij}=W_{ji}$ and $W_{ji}=W_{ij}^\top$ together give $W_{ij}=W_{ij}^\top$. That is $\tfrac{K(K-1)}{2}$ independent linear equations per unordered pair, and there are $\binom N2$ pairs. The remaining requirements are conic inequalities, which can only shrink the set further. $\square$

Proposition 6.6 is the quantitative form of the claim that the interaction shape is a restriction rather than a normalization. A generic Gaussian recognition law on $N$ agents does not have it, and no change of block-diagonal coordinates creates it, because the symmetry conditions are not congruence invariant. Section~\ref{sec:gauss-hypothesis} shows that generative constructions of the kind used in the preceding part do not supply it either.

\section{The interaction form is a hypothesis and not a consequence}
\label{sec:gauss-hypothesis}

\subsection{Trivialized coordinates and flat comparison}

The state latents carry frames. A residual between two agents is not $k_i-k_j$, which compares vectors in different fibers, but $k_i-\Omega_{ij}k_j$, where $\Omega_{ij}$ is the declared transport along the edge. The interaction shape of Definition 6.4 involves plain differences, so it can only be a statement in coordinates in which the transports have been removed. When that removal is possible is settled exactly by the following.

\paragraph{Proposition 6.7 (simultaneous trivialization is exactly coboundary).}
Let $G=(V,E)$ be a connected graph and let $\{\Omega_{ij}\}_{(i,j)\in E}\subset\GL^{+}(K)$ satisfy $\Omega_{ji}=\Omega_{ij}^{-1}$. There exist frames $\{U_i\}\subset\GL^{+}(K)$ with $U_i^{-1}\Omega_{ij}U_j=I$ for every edge if and only if $\Omega_{ij}=U_iU_j^{-1}$ for those frames, and this holds if and only if the ordered product of transports around every cycle of $G$ is the identity. When it holds, the frames are unique up to a common right factor $U_i\mapsto U_ig$ with $g\in\GL^{+}(K)$ independent of $i$. \status{ESTABLISHED}

\paragraph{Proof.}
The condition $U_i^{-1}\Omega_{ij}U_j=I$ is literally $\Omega_{ij}=U_iU_j^{-1}$, which gives the first equivalence. For the second, suppose such frames exist and let $i_0,i_1,\dots,i_r=i_0$ be a cycle. Then $\prod_{t}\Omega_{i_ti_{t+1}}=\prod_t U_{i_t}U_{i_{t+1}}^{-1}=U_{i_0}U_{i_0}^{-1}=I$ by telescoping. Conversely, assume every cycle product is the identity, fix a spanning tree and a root $r$, set $U_r=I$, and define $U_i$ as the ordered product of transports along the tree path from $r$ to $i$. The relation $\Omega_{ij}=U_iU_j^{-1}$ then holds on every tree edge by construction, and on a non-tree edge it holds precisely because the cycle formed by that edge and the tree path between its endpoints has trivial product. For uniqueness, if $U_iU_j^{-1}=U_i'U_j'^{-1}$ on every edge then $U_i'^{-1}U_i=U_j'^{-1}U_j$ for adjacent $i,j$, so this quantity is constant on a connected graph. $\square$

\paragraph{Hypothesis 6.8 (flat comparison).}
The declared transports are a coboundary, $\Omega_{ij}=U_iU_j^{-1}$, so that the trivialization $z_i=U_i^{-1}\mu_i$ of the means, and the same map $z_i=U_i^{-1}k_i$ applied to the latents themselves, turn every residual into a plain difference,
\begin{equation}
\mu_i-\Omega_{ij}\mu_j=U_i\left(z_i-z_j\right).
\label{eq:gauss-flat-trivialization}
\end{equation}
This is a choice, not a derivation. By Proposition 6.7 it excludes exactly the configurations with nontrivial holonomy around some cycle, which is the independently declared edge link of Regime II. It is used in two places: it is what makes the interaction shape of Definition 6.4 a statement about a fixed coordinate system at all, and it is what the closure of the family under aggregation requires. The mechanism of the second use is visible already: merging two agents identifies $z_i=z_j=z$, and the internal edge energy $(z_i-\widetilde\Omega_{ij}z_j)^\top W_{ij}(z_i-\widetilde\Omega_{ij}z_j)$ collapses to $z^\top(I-\widetilde\Omega_{ij})^\top W_{ij}(I-\widetilde\Omega_{ij})z$, where $\widetilde\Omega_{ij}$ is the trivialized transport. This residue vanishes for every $z$ exactly when $\widetilde\Omega_{ij}=I$, that is exactly under flat comparison. The coarse-graining chapter carries out the general computation and states what survives without the hypothesis. \status{HYPOTHESIS}

Two remarks fix the status of Equation~\eqref{eq:gauss-flat-trivialization}. First, the change of variables $y=Uz$ with $U=\operatorname{blockdiag}(U_1,\dots,U_N)$ acts on precisions by congruence, $\Lambda^{(y)}=U^{-\top}\Lambda^{(z)}U^{-1}$, which preserves inertia by Sylvester's law \citep[Ch.~4]{HornJohnson2013} but does not preserve membership in $\mathcal I(V)$, since the transformed off-diagonal block $U_i^{-\top}\Lambda^{(z)}_{ij}U_j^{-1}$ need not be symmetric when $U_i\neq U_j$. Membership in the interaction family is therefore a frame-dependent assertion. Second, that frame dependence is exactly why the invariants of the renormalization chapters are stated as generalized eigenvalues of the pair $(L,\Lambda)$ rather than as eigenvalues of either matrix alone: the pair transforms by the same congruence and the generalized spectrum is unchanged. \status{ESTABLISHED}

\subsection{What a directed linear-Gaussian model actually contributes}

Consider the state sector of a directed linear-Gaussian model of the kind used in the preceding part: each non-root agent $i$ has a single parent $j=\operatorname{pa}(i)$ and receiving kernel
\begin{equation}
k_i\given k_j\sim\mathcal N\left(G_i\Omega_{ij}k_j,\Lambda_i^{-1}\right),
\qquad
\Lambda_i\in\Sym_{++}^{K},
\qquad
G_i\in\GL(K),
\label{eq:gauss-receiving-kernel}
\end{equation}
and each root $r$ carries a proper prior $\mathcal N(0,\Lambda_r^{-1})$, a nonzero prior mean shifting only the information vector $h$ and not the precision. The gain is written $G_i$ rather than $A_i$ because the latter symbol is reserved for the self terms of Definition 6.4. A single subscript on $\Lambda$ denotes the innovation precision of one receiving kernel; a double subscript denotes a block of the assembled precision.

\paragraph{Proposition 6.9 (assembled precision blocks).}
Write $T_i:=G_i\Omega_{i\operatorname{pa}(i)}$ and $\operatorname{ch}(i)=\{\ell:\operatorname{pa}(\ell)=i\}$. The joint precision of the model above has blocks
\begin{equation}
\Lambda_{i\operatorname{pa}(i)}=-\Lambda_iT_i,
\qquad
\Lambda_{ii}=\Lambda_i+\sum_{\ell\in\operatorname{ch}(i)}T_\ell^\top\Lambda_\ell T_\ell,
\label{eq:gauss-directed-blocks}
\end{equation}
and $\Lambda_{ij}=0$ whenever $i$ and $j$ are not adjacent in the parent graph. \status{ESTABLISHED}

\paragraph{Proof.}
The negative log joint density is $\tfrac12\sum_i(k_i-T_ik_{\operatorname{pa}(i)})^\top\Lambda_i(k_i-T_ik_{\operatorname{pa}(i)})$ plus terms not depending on $k$, with the parent term absent at a root. Expanding one summand gives $\tfrac12k_i^\top\Lambda_ik_i-k_i^\top\Lambda_iT_ik_{\operatorname{pa}(i)}+\tfrac12k_{\operatorname{pa}(i)}^\top T_i^\top\Lambda_iT_ik_{\operatorname{pa}(i)}$. The precision is the Hessian of this function, and collecting the three contributions of each factor gives Equation~\eqref{eq:gauss-directed-blocks}. $\square$

By Equation~\eqref{eq:gauss-directed-blocks}, membership in the interaction family requires the matrix
\begin{equation}
W_{i\operatorname{pa}(i)}:=\Lambda_iT_i=\Lambda_iG_i\Omega_{i\operatorname{pa}(i)}
\label{eq:gauss-candidate-weight}
\end{equation}
to be symmetric and positive semidefinite. Neither property is automatic, and both fail on explicit parameter values.

Two sources of failure must be kept apart, and a convention does it. Under Hypothesis 6.8 the trivialized residual is $z_i-\widetilde T_iz_j$ with
\begin{equation}
\widetilde T_i:=U_i^{-1}T_iU_j=U_i^{-1}G_iU_i,
\qquad
\widetilde\Lambda_i:=U_i^\top\Lambda_iU_i,
\label{eq:gauss-trivialized-transition}
\end{equation}
the second equality in the first display using $\Omega_{ij}=U_iU_j^{-1}$, so that Equation~\eqref{eq:gauss-directed-blocks} holds verbatim in trivialized coordinates with $(\Lambda_i,T_i)$ replaced by $(\widetilde\Lambda_i,\widetilde T_i)$. The transports have then already been removed and any remaining obstruction is attributable to the gain alone. For the rest of this section the tildes are dropped and $(\Lambda_i,T_i)$ denote the trivialized quantities.

\paragraph{Proposition 6.10 (two witnesses).}
There exist admissible parameters for which no interaction-family representation of the assembled precision exists.
\begin{equation}
\Lambda_i=\begin{psmallmatrix}2&1\\1&3\end{psmallmatrix}\succ0,
\qquad
T_i=\begin{psmallmatrix}0&1\\-1&0\end{psmallmatrix}\in\GL^{+}(2)
\qquad\Longrightarrow\qquad
-\Lambda_iT_i=\begin{psmallmatrix}1&-2\\3&-1\end{psmallmatrix},
\label{eq:gauss-asymmetry-witness}
\end{equation}
which differs from its own transpose by a matrix of Frobenius norm $5\sqrt2$, so the off-diagonal block is not symmetric and no $W$ exists at all. Separately, taking $K=1$ with $G_i=-1$ and $\Omega_{ij}=\Lambda_i=1$ gives a symmetric block with $W=-1<0$, so the sign requirement fails; the same construction with $G_i=-I$, $\Omega_{ij}=I$ and $\Lambda_i=I$ gives $W=-I\prec0$ at every $K$. \status{ESTABLISHED}

\paragraph{Proof.}
Direct computation. For the first, $\Lambda_iT_i=\begin{psmallmatrix}-1&2\\-3&1\end{psmallmatrix}$, so $-\Lambda_iT_i$ is as displayed and $-\Lambda_iT_i-(-\Lambda_iT_i)^\top=\begin{psmallmatrix}0&-5\\5&0\end{psmallmatrix}$ with Frobenius norm $\sqrt{50}$. Positive definiteness of $\Lambda_i$ follows from $\det=5>0$ and trace $5>0$; invertibility of $T_i$ from $\det T_i=1>0$. For the second, Equation~\eqref{eq:gauss-candidate-weight} evaluates to $\Lambda_iG_i\Omega_{ij}=-I$, and $-I$ is not positive semidefinite for any $K\geq1$. $\square$

Proposition 6.10 settles the status of the interaction shape. The gain matrix $G_i$ is a free generative parameter in the construction of the preceding part, not required to be symmetric, positive, or even close to the identity, and for the displayed values the assembled precision has no representation of the form of Equation~\eqref{eq:gauss-interaction-family}. Whatever else may be said for the interaction family, it is not a consequence of the general linear-Gaussian directed model.

\subsection{Exactly when the form does hold}

Having shown that the form fails, the obligation is to say when it holds. Two statements are available, one global and one local, and they answer different questions.

\paragraph{Proposition 6.11 (exact global condition).}
The assembled precision of Equation~\eqref{eq:gauss-directed-blocks} lies in $\mathcal I(V)$ with the edge assignment of Equation~\eqref{eq:gauss-candidate-weight} if and only if, for every non-root $i$,
\begin{equation}
W_i:=\Lambda_iT_i\in\PSD^{K},
\label{eq:gauss-global-cond-1}
\end{equation}
and, for every agent $i$,
\begin{equation}
A_i:=\Lambda_i-W_i+\sum_{\ell\in\operatorname{ch}(i)}\left(W_\ell\Lambda_\ell^{-1}W_\ell-W_\ell\right)\succeq0,
\label{eq:gauss-global-cond-2}
\end{equation}
with the term $-W_i$ omitted at a root. \status{ESTABLISHED}

\paragraph{Proof.}
Condition~\eqref{eq:gauss-global-cond-1} is exactly the requirement that the off-diagonal blocks have the form $-W$ with $W$ symmetric positive semidefinite. Given it, $T_i=\Lambda_i^{-1}W_i$ and $T_\ell^\top\Lambda_\ell T_\ell=W_\ell\Lambda_\ell^{-1}\Lambda_\ell\Lambda_\ell^{-1}W_\ell=W_\ell\Lambda_\ell^{-1}W_\ell$, using symmetry of $W_\ell$. The self term is then the residual $A_i=\Lambda_{ii}-\sum_{j\sim i}W_{ij}$, where the edges at $i$ are the one to its parent, of weight $W_i$, and one to each child $\ell$, of weight $W_\ell$. Substituting Equation~\eqref{eq:gauss-directed-blocks} gives Equation~\eqref{eq:gauss-global-cond-2}, and family membership requires exactly that this residual be positive semidefinite. $\square$

Condition~\eqref{eq:gauss-global-cond-2} is a balance at each node and permits cancellation: a parent deficit $\Lambda_i-W_i\prec0$ may be offset by child surpluses. That freedom disappears if one asks the stronger and more natural question, namely whether each generative factor separately lies in the family, so that the interaction structure can be read off the model factor by factor rather than after a global cancellation. The answer is sharp.

\paragraph{Proposition 6.12 (edge-local characterization).}
Consider a single receiving factor $j\to i$ in isolation, contributing $\Lambda_i$ to the $ii$ block, $T_i^\top\Lambda_iT_i$ to the $jj$ block, and $-\Lambda_iT_i$ to the $ij$ block. This factor lies in $\mathcal I(\{i,j\})$ if and only if
\begin{equation}
X:=\Lambda_i^{-1/2}\left(\Lambda_iT_i\right)\Lambda_i^{-1/2}
\text{ is an orthogonal projection,}
\label{eq:gauss-edge-local}
\end{equation}
equivalently if and only if $\Lambda_iT_i$ is symmetric and $T_i$ is idempotent, equivalently if and only if $T_i$ is a $\Lambda_i$-orthogonal projection. In that case the self terms are $A_i=\Lambda_i^{1/2}(I-X)\Lambda_i^{1/2}$ and $A_j=0$. The full-rank case $X=I$ is $T_i=I$, in which the factor is the pure agreement factor whose mean is the parent's own value in the trivialized coordinates and whose edge weight is $W=\Lambda_i$. \status{ESTABLISHED}

\paragraph{Proof.}
Membership requires three things: $W:=\Lambda_iT_i\succeq0$ and symmetric; the $i$-side residual $\Lambda_i-W\succeq0$; and the $j$-side residual $T_i^\top\Lambda_iT_i-W\succeq0$. Assume $W$ symmetric and set $X=\Lambda_i^{-1/2}W\Lambda_i^{-1/2}$, which is symmetric. Then $W\succeq0$ reads $X\succeq0$; $\Lambda_i-W\succeq0$ reads $X\preceq I$; and, using $T_i=\Lambda_i^{-1}W$ so that $T_i^\top\Lambda_iT_i=W\Lambda_i^{-1}W=\Lambda_i^{1/2}X^2\Lambda_i^{1/2}$, the third reads $X^2\succeq X$. Now $0\preceq X\preceq I$ implies $X^2\preceq X$, since $X$ and $X^2$ are simultaneously diagonalizable and $\lambda^2\leq\lambda$ for $\lambda\in[0,1]$. Combining, $X^2=X$, so $X$ is a symmetric idempotent, that is an orthogonal projection. Conversely any orthogonal projection satisfies all three conditions, and then $T_i=\Lambda_i^{-1/2}X\Lambda_i^{1/2}$ satisfies $T_i^2=T_i$ and $\Lambda_iT_i=\Lambda_i^{1/2}X\Lambda_i^{1/2}$ is symmetric positive semidefinite. The self terms follow by substitution, the $j$-side residual being $\Lambda_i^{1/2}(X^2-X)\Lambda_i^{1/2}=0$. $\square$

Proposition 6.12 says what the interaction hypothesis amounts to at the level of the generative model: it demands that every coupling factor be an agreement factor, possibly restricted to a subspace, and never a general linear prediction. The degenerate cases with $X$ a proper projection are agreement factors that constrain only the components of the residual inside $\operatorname{range}(X)$ and leave the complement to the self term. This is a strong structural demand and the counterexamples of Proposition 6.10 are the generic situation rather than pathologies: a numerical sweep over $4000$ gain matrices drawn from a standard normal at $K=3$ against a fixed random $\Lambda_i$ (seed $20260728$) produced zero admissible cases, while the four projection cases of ranks $0$ through $3$ passed with idempotency residual below $3\times10^{-15}$. That sweep is a corroboration of the proved statement and not evidence for it. \status{NUMERICAL}

\paragraph{Hypothesis 6.13 (interaction form).}
The renormalization half of this document acts on recognition precisions of the form of Equation~\eqref{eq:gauss-interaction-family}, in trivialized coordinates, and this is adopted as a declared restriction. It excludes every coupling whose trivialized gain is not a $\Lambda_i$-orthogonal projection, in particular every generic linear-Gaussian transition, by Propositions 6.10 and 6.12. It is used in the aggregation closure statement, in the definition of the flow, and in the fixed-point analysis; each of those results names it. The document does not claim that it is forced by anything, and Proposition 6.6 records that it is a measure-zero restriction on the space of Gaussian recognition laws. \status{HYPOTHESIS}

\section{Nonemptiness, and a note on conditioning}
\label{sec:gauss-nonempty}

A declared subfamily is worthless if it is empty or numerically degenerate, so both should be checked. Nonemptiness is a construction, not an experiment.

\paragraph{Proposition 6.14 (constructive nonemptiness).}
Under Hypothesis 6.8, let each edge of a graph on $V$ carry an agreement factor with residual $\mu_i-\Omega_{ij}\mu_j$ and link covariance $R_{ij}\in\Sym_{++}^{K}$, and let each agent carry a proper prior of precision $A_i\succeq0$ in the trivialized coordinates. Then the assembled trivialized energy is
\begin{equation}
E(z)=\tfrac12\sum_{i}z_i^\top A_iz_i
+\tfrac12\sum_{i<j}\left(z_i-z_j\right)^\top W_{ij}\left(z_i-z_j\right),
\qquad
W_{ij}=U_i^\top R_{ij}^{-1}U_i,
\label{eq:gauss-nonempty-construction}
\end{equation}
so the associated precision lies in $\mathcal I(V)$, with each $W_{ij}$ symmetric positive definite by construction. \status{ESTABLISHED}

\paragraph{Proof.}
By Equation~\eqref{eq:gauss-flat-trivialization} the residual is $U_i(z_i-z_j)$, so the quadratic form $(\mu_i-\Omega_{ij}\mu_j)^\top R_{ij}^{-1}(\mu_i-\Omega_{ij}\mu_j)$ equals $(z_i-z_j)^\top U_i^\top R_{ij}^{-1}U_i(z_i-z_j)$. The matrix $U_i^\top R_{ij}^{-1}U_i$ is symmetric and positive definite because $R_{ij}^{-1}$ is and $U_i$ is invertible. Comparing with Equation~\eqref{eq:gauss-interaction-energy} identifies the parameters. $\square$

Proposition 6.14 also makes the role of Hypothesis 6.8 concrete by contrast. Without trivialization the same agreement factor contributes the off-diagonal block $-R_{ij}^{-1}\Omega_{ij}$, which is symmetric only in special position; a random draw at $K=3$ (seed $20260728$) differed from its own transpose by a matrix of Frobenius norm $4.32$, while the corresponding trivialized weight was symmetric to $1.8\times10^{-16}$ with smallest eigenvalue $1.07\times10^{-2}$; the trivialized weight is the control, and it is the case Proposition 6.14 proves. Computation is not proof. \status{NUMERICAL}

Conditioning is a different question and admits only computational evidence. Drawing $N=6$ agents at $K=3$ with each unordered pair present independently with probability $0.6$, edge weights $W_{ij}=MM^\top+0.1I$ and self terms $A_i=MM^\top/K$ for independent standard normal $M$ (seed $20260728$, $200$ draws), every assembled $\Lambda$ was positive definite, with smallest observed eigenvalue $4.6\times10^{-4}$ and condition numbers of median $93.4$ and maximum $4.0\times10^{4}$. Two controls were run. Setting $A_i=0$ produced exactly $K=3$ zero eigenvalues with $\lVert L(\mathbf 1\otimes I_K)\rVert_F=2.9\times10^{-15}$, which is the consensus null space predicted by Proposition 6.5. Drawing $200$ unstructured positive definite matrices of the same size produced zero matrices with symmetric off-diagonal blocks, corroborating Proposition 6.6. Computation is not proof: positive definiteness here is proved by Proposition 6.5 and thinness by Proposition 6.6, and the measurements speak only to the conditioning of one sampling scheme, for which no proof is offered and none is claimed. \status{NUMERICAL}

\section{What is not settled}
\label{sec:gauss-open}

One question about the family itself is left open, and it is worth stating precisely because it is open in the algebraic-graph literature as well.

\paragraph{Open problem 6.15 (marginalization).}
Aggregation, treated in the coarse-graining chapter, merges agents. The complementary operation eliminates them by marginalization, replacing $\Lambda$ by the Schur complement $\Lambda_{RR}-\Lambda_{RE}\Lambda_{EE}^{-1}\Lambda_{ER}$ on a retained set $R$ with eliminated set $E$. In the scalar-weighted case this is Kron reduction and returns a Laplacian. The matrix-weighted case is not settled. What can be stated exactly is the obstruction for a single eliminated agent $e$: the new off-diagonal block between retained agents $i$ and $j$ is $-W_{ij}-W_{ie}\Lambda_{ee}^{-1}W_{ej}$, and since the block is determined there is no freedom in the edge assignment, so membership requires $W_{ie}\Lambda_{ee}^{-1}W_{ej}$ to be symmetric. That is the commutation condition $W_{ie}\Lambda_{ee}^{-1}W_{ej}=W_{ej}\Lambda_{ee}^{-1}W_{ie}$, which is automatic at $K=1$. \status{ESTABLISHED} That it fails generically above $K=1$ is measured rather than proved: over $200$ independent draws at $K=3$ (seed $20260728$) it held zero times, with relative asymmetry up to $1.53$, while the $K=1$ control held $200$ times out of $200$. Computation is not proof, and no proof of generic failure is offered. \status{NUMERICAL} What remains open is whether some structurally motivated subfamily, for instance one with simultaneously diagonalizable edge weights, is closed under marginalization, and whether the positive semidefinite row-sum condition survives when the symmetry condition is imposed. Settling it requires either such a closure proof or a demonstration that no subfamily larger than the commuting one works. \status{OPEN}

\section{Status register}

{\small
\begin{longtable}{@{}>{\raggedright\arraybackslash}p{0.12\linewidth}>{\raggedright\arraybackslash}p{0.29\linewidth}>{\raggedright\arraybackslash}p{0.16\linewidth}>{\raggedright\arraybackslash}p{0.33\linewidth}@{}}
\caption{Epistemic status of the results of this chapter. Every entry not marked \textsc{established} names what is owed.}
\label{tab:gauss-register}\\
\toprule
\textbf{Item} & \textbf{Statement} & \textbf{Status} & \textbf{What is owed} \\
\midrule
\endfirsthead
\multicolumn{4}{@{}l}{\emph{Status register for Chapter~\ref{ch:gaussian}, continued.}}\\
\toprule
\textbf{Item} & \textbf{Statement} & \textbf{Status} & \textbf{What is owed} \\
\midrule
\endhead
\bottomrule
\endlastfoot
Prop.~6.1 & Log normalizer, moments, convexity of the Gaussian information form & ESTABLISHED & Nothing; proved, with the regularity condition cited \\
Prop.~6.2 & Conditional mean and covariance; zero precision block equals conditional independence & ESTABLISHED & Nothing; proved and cited \\
Rem.~6.3 & A precision block is not a covariance value; marginal block is an inverse Schur complement & ESTABLISHED & Nothing; explicit witness given \\
Def.~6.4 & The interaction family $\mathcal I(V)$ & DEFINITION & Nothing is proved by a definition \\
Prop.~6.5 & Energy form, semidefiniteness, exact definiteness criterion, consensus kernel & ESTABLISHED & Nothing; proved \\
Prop.~6.6 & The family is Lebesgue-null in $\Sym^{NK}$ for $K\geq2$, $N\geq2$ & ESTABLISHED & Nothing; dimension count \\
Prop.~6.7 & Simultaneous trivialization holds exactly for coboundary transports & ESTABLISHED & Nothing; proved both directions \\
Hyp.~6.8 & Flat comparison & HYPOTHESIS & A declared choice; excludes nontrivial holonomy. Used by the aggregation closure, which must restate it \\
Prop.~6.9 & Assembled precision blocks of a directed linear-Gaussian model & ESTABLISHED & Nothing; proved \\
Prop.~6.10 & The interaction form fails: asymmetry and sign witnesses & ESTABLISHED & Nothing; arithmetic \\
Prop.~6.11 & Exact global condition for family membership & ESTABLISHED & Nothing; proved \\
Prop.~6.12 & Edge-local membership holds exactly for $\Lambda_i$-orthogonal projections & ESTABLISHED & Nothing; proved. The accompanying sweep is corroboration only \\
Hyp.~6.13 & The interaction form is adopted as a restriction & HYPOTHESIS & A declared choice; excludes generic linear transitions. Every later use must name it \\
Prop.~6.14 & The family is nonempty; explicit agreement-factor construction & ESTABLISHED & Nothing; proved \\
Prop.~6.14 & Measured contrast between the untrivialized and trivialized agreement weights & NUMERICAL & Computation only, seed $20260728$, the trivialized weight serving as control; the statement is proved by Prop.~6.14 \\
Sec.~6.5 & Conditioning of one sampling scheme & NUMERICAL & Computation only, with two controls; no conditioning bound is proved or claimed \\
Open~6.15 & Exact symmetry obstruction to one-agent marginalization & ESTABLISHED & Nothing; the reduced block is determined, so membership requires $W_{ie}\Lambda_{ee}^{-1}W_{ej}$ symmetric \\
Open~6.15 & That the commutation condition fails generically above $K=1$ & NUMERICAL & Computation only, seed $20260728$, with the $K=1$ case as control; no proof of generic failure is offered \\
Open~6.15 & Whether any subfamily is closed under marginalization & OPEN & A closure proof for a named subfamily, or a demonstration that none larger than the commuting case exists \\
\end{longtable}
}
